A fleet weighing the cost of collecting and combining evidence needs to know
not that ignoring heterogeneity is wrong, which is established, but how many
spares it changes. This section answers that for one population, first exactly
for an idealised case and then for the shortcuts a planner is likely to use.

\subsection{Setting and Exact Result}
\begin{assumption}
\label{as:mix}
A population expects $\lambda$ failures over the mission under a predictor that
ignores latent log-rate spread of standard deviation $\sigma$. The ignored
spread acts multiplicatively on the rate and preserves the mean:
\begin{equation}
Y\mid u\sim\mathrm{Poisson}\!\left(\lambda e^{u-\sigma^2/2}\right),\qquad
u\sim\mathcal N(0,\sigma^2).
\label{eq:mixture}
\end{equation}
\end{assumption}

Assumption~\ref{as:mix} is obtained from the model of Section~\ref{sec:model}
for a population at an unseen holder by collapsing the omitted latent terms into
one lognormal factor. The collapse is exact only if those terms are Gaussian on
the log-rate scale and act on the whole population alike, the cell dispersion
is absent or absorbed into $\sigma^2$, and parameter uncertainty is ignored.
The results below are therefore diagnostic. A plan for a real population is
computed from the predictive distribution of Section~\ref{sec:model}, as
Sections~\ref{sec:design}--\ref{sec:lighting} do.

\begin{proposition}[Exact delivered cover]
\label{prop:exact}
Under Assumption~\ref{as:mix} a stock $q$ is not exceeded with probability
\begin{equation}
C(q)=\int F_{\mathrm{Pois}}\!\left(q;\lambda e^{u-\sigma^2/2}\right)
\varphi_\sigma(u)\,\mathrm du ,
\label{eq:cover}
\end{equation}
where $\varphi_\sigma$ is the $\mathcal N(0,\sigma^2)$ density, and the smallest
stock with cover at least $p$ is $q^\ast_p=\min\{q:C(q)\ge p\}$.
\end{proposition}
\begin{IEEEproof}
Condition on $u$, apply the Poisson distribution function and integrate
against the law of $u$; $q$ does not depend on $u$.
\end{IEEEproof}

The integral is one-dimensional and cheap. Its integrand is a sigmoid in $u$
whose transition narrows as $\lambda$ grows, and Gauss--Hermite quadrature with
60 nodes misplaces it by up to \quadGH{} percentage points; a composite Simpson
rule on $[-10\sigma,10\sigma]$ with \quadNodes{} points is stable to
$\quadGrid$ and is used throughout.

\subsection{Shortcut Plans}
Three shortcuts are within reach of a planner who has an estimate of $\sigma$.
The first ignores it and embarks the Poisson quantile $q^{\mathrm P}_p$. The
second scales the Poisson safety margin,
\begin{equation}
q^{\kappa}_p=\big\lceil\lambda+(q^{\mathrm P}_p-\lambda)\sqrt{1+\kappa^2}\,\big\rceil,
\qquad \kappa=\sigma\sqrt\lambda ,
\label{eq:kappa}
\end{equation}
which is what a normal approximation gives: with $Y\mid u$ approximately normal
and $u$ small, $\mathrm{Var}(Y)\approx\lambda(1+\kappa^2)$. The same group
appears in capacity sizing under an uncertain arrival rate, where the
coefficient of variation of the rate is compared with the inverse square root
of the load~\cite{bassamboo2010,mathijsen2018}; moment-based safety stocks of
this kind are known to miss their target when the demand distribution is
skewed~\cite{eppen1988}. The third replaces the lognormal mixture by the
negative binomial with the same mean and variance, shape
$1/(e^{\sigma^2}-1)$, the gamma--Poisson form common in spare-parts
models~\cite{aronis2004,sherbrooke2004}.

Table~\ref{tab:single} compares the three with $q^\ast_p$ over
\spLamN{} values of $\lambda\in[\spLamLo,\spLamHi]$, ten values of
$\sigma\in[\spSigLo,\spSigHi]$ and three levels, \spPts{} points in all, a
range chosen to cover a single reference on a single ship. These are maxima
observed on a grid, not analytical bounds.

Ignoring $\sigma$ is costly: at the 99\,\% level the Poisson plan falls short
of its target by up to \spPoisShort{} percentage points, and by
\spPoisShortOp{} points even for $\lambda\le 20$ and $\sigma\le0.6$. The
margin factor removes most of that loss but not all of it. It still falls
\spKapShort{} points short, \spKapShortOp{} in the restricted range, and
matches the exact plan at only \spKapEqual\,\% of the grid points, because
both the normal approximation and the linearisation in $u$ fail for small
$\lambda$ and large $\kappa$. The moment-matched negative binomial is the
useful shortcut. At the 99\,\% level it equals the exact plan at
\spNbEqual\,\% of the points and is never more than \spNbShort{} points short,
because the gamma mixture has a lighter far tail than the lognormal, which
costs at most \spNbUnits{} units at the largest $\lambda\sigma$ explored. At 90
and 95\,\% the error changes sign and it over-provisions by up to
\spNbOver{} units.

\begin{table}[!t]
\caption{Shortcut plans against the exact plan of
Proposition~\ref{prop:exact} (\textsc{analytical}; grid maxima, not bounds):
largest shortfall of delivered cover below the target, in percentage points,
and the range of the difference from $q^\ast_p$ in units}
\label{tab:single}
\centering
\footnotesize
\setlength{\tabcolsep}{3pt}
\begin{tabular}{@{}lcccccc@{}}
\toprule
& \multicolumn{2}{c}{90\,\%} & \multicolumn{2}{c}{95\,\%} & \multicolumn{2}{c}{99\,\%} \\
\cmidrule(lr){2-3}\cmidrule(lr){4-5}\cmidrule(l){6-7}
Plan & short & units & short & units & short & units \\
\midrule
\tabinput{tables/tab_single.tex}
\bottomrule
\end{tabular}
\end{table}

\subsection{Use of the Result}
The recommendation is narrow. When a single reference is planned on its own,
the negative binomial with the fitted total spread is an adequate shortcut at
high cover levels, and the margin factor should not be used to size stock. Two
things the shortcuts do not handle are the subject of the rest of the paper:
spares are held as a vector over several references, and the spread $\sigma$
is itself estimated, with an uncertainty the plan has to carry.
